% Options for packages loaded elsewhere
\PassOptionsToPackage{unicode}{hyperref}
\PassOptionsToPackage{hyphens}{url}
%
\documentclass[
  brazilian,
]{article}
\usepackage{lmodern}
\usepackage{amssymb,amsmath}
\usepackage{ifxetex,ifluatex}
\ifnum 0\ifxetex 1\fi\ifluatex 1\fi=0 % if pdftex
  \usepackage[T1]{fontenc}
  \usepackage[utf8]{inputenc}
  \usepackage{textcomp} % provide euro and other symbols
\else % if luatex or xetex
  \usepackage{unicode-math}
  \defaultfontfeatures{Scale=MatchLowercase}
  \defaultfontfeatures[\rmfamily]{Ligatures=TeX,Scale=1}
\fi
% Use upquote if available, for straight quotes in verbatim environments
\IfFileExists{upquote.sty}{\usepackage{upquote}}{}
\IfFileExists{microtype.sty}{% use microtype if available
  \usepackage[]{microtype}
  \UseMicrotypeSet[protrusion]{basicmath} % disable protrusion for tt fonts
}{}
\makeatletter
\@ifundefined{KOMAClassName}{% if non-KOMA class
  \IfFileExists{parskip.sty}{%
    \usepackage{parskip}
  }{% else
    \setlength{\parindent}{0pt}
    \setlength{\parskip}{6pt plus 2pt minus 1pt}}
}{% if KOMA class
  \KOMAoptions{parskip=half}}
\makeatother
\usepackage{xcolor}
\IfFileExists{xurl.sty}{\usepackage{xurl}}{} % add URL line breaks if available
\IfFileExists{bookmark.sty}{\usepackage{bookmark}}{\usepackage{hyperref}}
\hypersetup{
  pdftitle={Parte III: Processamento de Linguagem Natural para Linguistas},
  pdfauthor={CiberExt 26-29 · FEELT38103 · Universidade Federal de Uberlândia},
  pdflang={pt-BR},
  hidelinks,
  pdfcreator={LaTeX via pandoc}}
\urlstyle{same} % disable monospaced font for URLs
\setlength{\emergencystretch}{3em} % prevent overfull lines
\providecommand{\tightlist}{%
  \setlength{\itemsep}{0pt}\setlength{\parskip}{0pt}}
\setcounter{secnumdepth}{-\maxdimen} % remove section numbering
\ifxetex
  % Load polyglossia as late as possible: uses bidi with RTL langages (e.g. Hebrew, Arabic)
  \usepackage{polyglossia}
  \setmainlanguage[]{brazilian}
\else
  \usepackage[shorthands=off,main=brazilian]{babel}
\fi

\title{Parte III: Processamento de Linguagem Natural para Linguistas}
\usepackage{etoolbox}
\makeatletter
\providecommand{\subtitle}[1]{% add subtitle to \maketitle
  \apptocmd{\@title}{\par {\large #1 \par}}{}{}
}
\makeatother
\subtitle{Unidade 1 - Exercícios}
\author{CiberExt 26-29 · FEELT38103 · Universidade Federal de
Uberlândia}
\date{Agosto de 2026}

\begin{document}
\maketitle

\hypertarget{exercuxedcios-pruxe1ticos-processamento-de-linguagem-natural-para-linguistas}{%
\section{Exercícios Práticos: Processamento de Linguagem Natural para
Linguistas}\label{exercuxedcios-pruxe1ticos-processamento-de-linguagem-natural-para-linguistas}}

Este documento contém exercícios práticos baseados na unidade
``Processando Diversos Idiomas''.

\hypertarget{muxf3dulo-1-processando-diversos-idiomas}{%
\subsection{Módulo 1: Processando Diversos
Idiomas}\label{muxf3dulo-1-processando-diversos-idiomas}}

\textbf{Exercício 1:} Historicamente, as ferramentas de Processamento de
Linguagem Natural (PNL) foram desenvolvidas com forte foco no idioma
inglês. Explique dois desafios estruturais significativos encontrados ao
tentar aplicar essas mesmas ferramentas convencionais para outros
idiomas, como o finlandês, turco ou o japonês.

\textbf{Exercício 2:} Utilizando a biblioteca Stanza, escreva o código
Python necessário para fazer o download do modelo de linguagem para o
idioma finlandês (código de idioma
\texttt{\textquotesingle{}fi\textquotesingle{}}). Em seguida, inicialize
um pipeline (\emph{Pipeline}) que inclua exclusivamente os processadores
de tokenização e etiquetagem de classes gramaticais (POS).

\textbf{Exercício 3:} Após processar um texto no Stanza, a saída gerada
para uma sentença não é uma lista comum do Python, mas um objeto
\emph{Sentence}. Qual método nativo da biblioteca você pode utilizar
para converter as anotações linguísticas desse objeto em uma lista de
dicionários Python, facilitando a interação com as anaves/valores?

\textbf{Exercício 4:} Qual é a maneira mais eficiente, em termos de
custo computacional e estruturação de código, para processar múltiplos
textos (documentos) de uma só vez utilizando o Stanza? Descreva os
passos técnicos utilizando a sintaxe de compreensão de listas
(\emph{list comprehension}).

\textbf{Exercício 5:} A biblioteca \texttt{spacy-stanza} é uma excelente
ferramenta para integrar os modelos pesados do Stanza dentro do
ecossistema familiar do spaCy. Apesar dessa vantagem, o texto menciona
uma ``grande limitação'' para o seu uso. Qual é essa limitação? Dê um
exemplo de idioma que ilustra esse problema.

\begin{center}\rule{0.5\linewidth}{0.5pt}\end{center}

\end{document}
